\scanpage{210}\mspage{155}
ou encore \unsure{(compte|compte} tenu de l'équivalence \inserted{induite par l'inclusion}
\[
 (\underline{\mathrm{Liss}}\,S)(W_S^d)^{-1}\xrightarrow{\approx}(\underline{\underline{\mathrm{Cat}}}_{/S})(W_S^d)^{-1}
\]
un foncteur, encore noté $\mathrm{mult}^{\mathrm{hot}}_{\deleted{\unsure{X/S}}X/S}$
\[
 \mathrm{mult}^{\mathrm{hot}}_{X/S}:\mathrm{HOT}_{\unsure{W|w}}(S)\unsure{-|\longrightarrow}\mathrm{HOT}_W(S).
\]
\deleted{Sa valeur} Ce foncteur satisfait donc la relation
\[
 \mathrm{mult}^{\mathrm{hot}}_{X/S}(\mathrm{hot}_{S,W}(Y))\overset{\mathrm{iso.can}}{\simeq}\mathrm{hot}_{S,W}(X\times_S Y)
\]
pour $Y$ \uline{lisse sur} $S$,
\margin{$\unsure{\mathrm{Z}|2}$}
et pour $Y$ quelconque dans $\mathrm{Cat}_{/S}$, \unsure{à|''}
\[
 \mathrm{mult}^{(\mathrm{hot})}_{X/S}(\mathrm{hot}_{S,W}(Y))\overset{\mathrm{iso\,can}}{\simeq}\mathrm{hot}_{S,W}(X\times_S\Phi_S(Y)).
\]
\deleted{Mais on se gardera de confon\unsure{dre}} Quand $X$ est lui-même lisse sur $S$, ce foncteur \deleted{\unsure{est donc}} n'est autre que la multiplication par $\mathrm{hot}_{S,W}(X)$ :
\[
 \mathrm{mult}^{(\mathrm{hot})}_{X/S}=\mathrm{mult}_{\mathrm{hot}_{S,W}(X)}
 \qquad\text{si }X\text{ \uline{lisse sur }}S.
\]
\inserted{Mais} \unsure{Sa|sa} valeur pour \deleted{$\unsure{Y}$ \unsure{lisse sur $S$}} $\mathrm{hot}_{S,W}(Y)$\unsure{,|$\ast$} quand $Y$ n'est \uline{pas} lisse sur $S$, ne doit pas être confondue avec $\mathrm{hot}_{S,W}(X\times_S Y)$.
\margin{$\mathrm{Z}$}

Ainsi, si $X$ et $Y$ sont dans $\mathrm{Cat}_{/S}$, considérons :
\[
 X\times_S Y\qquad\text{et}\qquad
 \tikz[baseline=(h.base)]{\node[ellipse,inner xsep=4pt,inner ysep=2pt] (h) {$\mathrm{hot}_{S,W}(X\times_S Y).$}; \draw (h.18) to[out=155,in=5] (h.north) to[out=180,in=90] (h.west) to[out=-90,in=180] (h.south) to[out=0,in=-90] (h.east) to[out=90,in=0] (h.6);}
\]
a) Si ni $X$, ni $Y$ ne sont supposés lisses sur $S$, \unsure{ce dernier|ça} ``\unsure{n'a} aucune interprétation \unsure{naturelle} en termes d'opérations \deleted{entre catégories} \inserted{sur \deleted{des} catégories} \deleted{dans des} $\mathrm{HOT}_W(\unsure{S|s})$''.\textsuperscript{\unsure{$\ast$}}
\margin{\unsure{Sauf} dans le cas où l'une des deux est dans $\underline{\unsure{\mathrm{Fibhot}|\mathrm{FibHot}}}\,S$, \unsure{comme on vient de voir}.}
% slanted
% archive sheet number: 209
